%% sec_results.tex
%% ─────────────────────────────────────────────────────────────────────────────
%% Section 6 — Results
%%
%% NOTE: Sections marked \TODO{} require measured data not yet available.
%%       Placeholder tables/figures include the structure and known values.
%% ─────────────────────────────────────────────────────────────────────────────

\section{Results}\label{sec:results}

\subsection{Program memory utilisation}\label{sec:res_pm}

Table~\ref{tab:pm_results} summarises the program memory reported by the AIE
compiler for each kernel after applying all optimisations described in
Section~\ref{sec:implementation}.  All five kernels fit within the 16\KB
limit.  The critical path was \texttt{K1}: the lazy-evaluation optimisation
that moved \texttt{FFV1P0\_3} from \texttt{K1} to \texttt{K3} reduced its
footprint from ${\sim}17.8\KB$ to ${\sim}15.5\KB$ (saving ${\sim}2.7\KB$,
$\Delta = -15\%$).

\begin{table}[h]
\caption{Program memory utilisation per kernel.  The 16\KB hardware limit is
16\,384 bytes.  Usage percentages are relative to that limit.}
\label{tab:pm_results}
\begin{tabular}{@{}lrrrl@{}}
\toprule
Kernel & PM (bytes) & Usage (\%) & Headroom (bytes) & Status \\
\midrule
K1 — WFGEN          & ${\sim}15{,}500$ & ${\sim}94.6$ & ${\sim}884$ & \textcolor{orange}{near-full} \\
K2a — FF\_DIAG\_A   & ${\sim}13{,}000$ & ${\sim}79.3$ & ${\sim}3{,}384$ & OK \\
K2b — FF\_DIAG\_B   & ${\sim}12{,}600$ & ${\sim}76.9$ & ${\sim}3{,}784$ & OK \\
K3 — VVV\_AMP       & ${\sim}9{,}600$  & ${\sim}58.6$ & ${\sim}6{,}784$ & OK \\
K4 — VVVV\_COLOR    & ${\sim}8{,}400$  & ${\sim}51.3$ & ${\sim}7{,}984$ & OK \\
\bottomrule
\end{tabular}
\end{table}

\subsection{Physics equivalence and float32 precision}\label{sec:res_prec}

The cascade pipeline was validated against the \mgfive double-precision
(\texttt{float64}) reference at multiple levels:
\begin{itemize}
  \item \textbf{Wavefunction level}: external wavefunctions $w_0$--$w_4$
    computed by \texttt{K1} agree with the \mgfive reference to within
    float32 rounding.
  \item \textbf{Diagram level}: amplitudes D1--D16 (compared at the JAMP
    accumulation stage in \texttt{K2a}, \texttt{K2b}, \texttt{K3}, \texttt{K4})
    have component-wise agreement with the reference.
  \item \textbf{JAMP level}: the full \texttt{JAMP[0..5]} vector at the output
    of \texttt{K3} matches the reference.
  \item \textbf{$|M|^2$ level}: final \ME outputs are compared against the
    \mgfive double-precision result.
\end{itemize}

\TODO{Insert precision scatter plot and relative-error histogram for 1000-PSP
study.  Include mean and maximum relative error, and comparison with CUDACPP
float32 mode.}

The float32 relative precision is ${\sim}10^{-7}$, far below the physical
uncertainties in LO calculations (${\sim}1\%$).  Systematic errors (rather
than random) arise from the float32 accumulation order, as expected.

\subsection{Single-pipeline latency}\label{sec:res_latency}

\TODO{Insert measured single-pipeline latency from hardware run on VCK190.
Expected: ${\sim}81\,\mu$s per PSP based on simulation and cycle estimates.}

The single-pipeline latency is the time from PSP injection at \texttt{plio\_in}
to ME$^2$ read-back at \texttt{plio\_out}, measured with XRT profiling.
This corresponds to \textbf{\TODO{X}~$\mu$s per PSP}, giving a single-pipeline
throughput of $1/T_\text{pipe}\ \text{PSP/s}$.

\subsection{PLIO bandwidth budget}\label{sec:res_plio}

The architecture is overwhelmingly \emph{compute-bound}, not I/O-bound.  Each
pipeline consumes one PSP of 80 bytes and emits one ME$^2$ result of 4 bytes
($+8$ bytes packet overhead).  Per-trunk PLIO capacity at 500\,MHz and 32-bit
width is 2\,GB/s.  With 8 pipelines sharing one PLIO trunk:
\begin{align}
  \text{PLIO utilisation (input)} &= \frac{8\times 80\,\text{B}/T_\text{pipe}}{2\,\text{GB/s}}
    \approx 0.3\%, \notag \\
  \text{PLIO utilisation (output)} &\ll 0.1\%.
\end{align}
The PLIO interface could accept data ${\sim}300\times$ faster than the
compute pipeline processes it.  PSP injection time per packet (21 beats at
500\,MHz) is ${\sim}42\,\text{ns}$, compared to ${\sim}81\,\mu\text{s}$ compute
time---a ratio of ${\sim}1{,}900\times$.  This confirms that all
performance optimisation effort should target kernel compute latency.

\subsection{80-pipeline aggregate throughput and energy}\label{sec:res_throughput}

With 80 independent, non-contending pipelines:
\begin{equation}
  \text{Throughput}_{\text{80-pipe}}
  = \frac{80}{T_\text{pipe}}\ \text{PSP/s}
  \approx \mathbf{\TODO{${\sim}1\,\text{MPSP/s}$}},
  \label{eq:throughput}
\end{equation}
where linear scaling is justified because each pipeline is physically
independent (separate cascade chains, separate packet-switched I/O, no shared
memory or synchronisation).

\TODO{Confirm with hardware measurement or emulation.  Provide energy estimate
at 45W board power: $E = 45\,\text{W} / T_{\text{throughput}}$ nJ/PSP.}

\subsection{Platform comparison}\label{sec:res_comparison}

Table~\ref{tab:comparison} places the AIE result in context of
existing acceleration efforts.  GPU and CPU numbers are sourced from
published CUDACPP benchmarks~\cite{Valassi:2025xfn}.

\begin{table}[h]
\caption{Performance and energy comparison for \ggttg matrix element
evaluation across platforms.  \TODO{Fill AIE measured values once hardware
benchmark is complete.  GPU/CPU numbers from~\cite{Valassi:2025xfn}.}}
\label{tab:comparison}
\begin{tabular}{@{}lllll@{}}
\toprule
Platform & Cores / tiles & Throughput & Power & Energy/PSP \\
\midrule
AMD VCK190 — \aie (80 pipes) &
  400 AIE tiles &
  \TODO{${\sim}1$\,MPSP/s} &
  ${\sim}45$\,W &
  \TODO{${\sim}46$\,nJ} \\
AMD VCK190 — 1 pipe (measured) &
  5 AIE tiles &
  ${\sim}12.3$\,kPSP/s &
  ${\sim}45$\,W &
  ${\sim}3.6\,\mu$J \\
NVIDIA V100 (CUDACPP float32) &
  5\,120 CUDA cores &
  \TODO{est.} &
  300\,W &
  \TODO{est.~${\sim}60$\,nJ} \\
Intel Xeon (single core, MG5aMC) &
  1 core &
  ${\sim}10$\,kPSP/s &
  150\,W (TDP) &
  ${\sim}10\,\mu$J \\
Intel Xeon (16 cores, MG5aMC) &
  16 cores &
  ${\sim}100$\,kPSP/s &
  150\,W (TDP) &
  ${\sim}1.5\,\mu$J \\
\bottomrule
\end{tabular}
\end{table}

\subsection{Resource utilisation summary}\label{sec:res_resources}

The design achieves \textbf{100\% \aie tile utilisation} (400 of 400 tiles
occupied by the 80 five-tile pipelines).  PL resource consumption is modest:
10$\times$ MM2S and 10$\times$ S2MM HLS kernels each consume $\ll 1\%$ of
available LUTs.

\TODO{Insert PL LUT/DSP utilisation table from Vivado/Vitis reports.}
